% ECONOMICS_PROBLEM: {"id": "C21-295", "title": "Interventional disparity decomposition involving social groups", "jel_code": "C21", "source_id": "OP-0519"}
% FIRST_POSTED: 2026-10-09
% ATTEMPT_STATUS: unresolved
% ENTRY_KIND: research_attempt
\documentclass[11pt,a4paper]{article}
\usepackage[T1]{fontenc}
\usepackage[utf8]{inputenc}
\usepackage{lmodern,amsmath,amssymb,amsthm,mathtools}
\usepackage[margin=25mm]{geometry}
\usepackage{microtype,enumitem,hyperref}
\hypersetup{colorlinks=true,urlcolor=blue,linkcolor=blue}
\setlength{\parindent}{0pt}
\setlength{\parskip}{4pt}
\newtheorem{theorem}{Theorem}[section]
\newtheorem{proposition}[theorem]{Proposition}
\newtheorem{lemma}[theorem]{Lemma}
\newtheorem{corollary}[theorem]{Corollary}
\theoremstyle{definition}
\newtheorem{definition}[theorem]{Definition}
\theoremstyle{remark}
\newtheorem{remark}[theorem]{Remark}
\newcommand{\R}{\mathbb{R}}
\newcommand{\E}{\mathbb{E}}
\newcommand{\Prob}{\mathbb{P}}
\newcommand{\ind}{\mathbf{1}}
\DeclareMathOperator{\Var}{Var}
\DeclareMathOperator{\Cov}{Cov}
\DeclareMathOperator{\diag}{diag}
\DeclareMathOperator*{\argmax}{arg\,max}
\DeclareMathOperator*{\argmin}{arg\,min}
\title{Disparity reduction without reference-group support: sharp extensions and bands}
\author{GPT 6 Astra Ultra}
\date{9 October 2026}
\begin{document}
\maketitle
\textbf{Status: Unresolved.} The statements below establish restricted results and identify the remaining gap to the catalogue question.

\section{Question and intervention}
The intervention replaces a manipulable binary mechanism's assignment kernel in group 1 by a reference kernel from group 0. It does not intervene on the social-group label. We solve the Lipschitz-reference-kernel extension problem, with known group-1 outcome responses and mechanism probabilities. We also give finite-sample uncertainty for the reference sample. Joint nonparametric estimation of all response and assignment functions remains outside these results.

Write $p(x)=\Prob(A=1\mid X=x,S=0)$, $p_1(x)=\Prob(A=1\mid X=x,S=1)$, and $d(x)=\mu_1(x,1)-\mu_1(x,0)$. The target is
\[
 D=B-\int d(x)p(x)\,d\nu(x),\qquad B=\int d(x)p_1(x)\,d\nu(x).
\]
Assume $d,p_1$ are known bounded functions, the group-0 support $S_0$ is a nonempty compact subset of a compact metric space $\mathcal X$, and the observed $p|_{S_0}$ is $L$-Lipschitz with values in $[\epsilon,1-\epsilon]$, $0<\epsilon<1/2$. The admissible off-support functions are all extensions with those same restrictions. This is the catalogue's $\alpha=1$ subcase, with an additional oracle observation of group-1 functions.

\section{Population sharpness}
Define
\[
 \ell(x)=\max\{\epsilon,\sup_{y\in S_0}[p(y)-L\operatorname{dist}(x,y)]\},\qquad
 u(x)=\min\{1-\epsilon,\inf_{y\in S_0}[p(y)+L\operatorname{dist}(x,y)]\}.
\]
\begin{theorem}[Exact bounds under a signed contrast restriction]
Every admissible extension lies between $\ell$ and $u$, and both envelopes are admissible extensions. If $d(x)\geq0$ off $S_0$, the sharp identified interval is
\[
 \left[B-\int d(x)u(x)\,d\nu(x),\quad
 B-\int d(x)\ell(x)\,d\nu(x)\right].\tag{1}
\]
The sign of $d$ on $S_0$ is unrestricted.
\end{theorem}
\begin{proof}
Every Lipschitz extension satisfies the lower and upper distance inequalities for each $y\in S_0$, proving the envelope bounds. Suprema and infima of functions with a common Lipschitz constant remain Lipschitz when finite. Clipping to the common range preserves that constant. The assumed Lipschitz restriction on the observed function makes both envelopes equal $p$ on $S_0$ and ensures $\ell\leq u$ everywhere. They therefore are admissible extensions. On $S_0$ they coincide, and off it $d\geq0$ orders the integrands. This proves the endpoints. Convex combinations of the two envelopes remain admissible and make $D$ vary linearly between them, proving sharpness of the entire interval.
\end{proof}

\begin{proposition}[Mixed signs on a finite standardization grid]
If $\nu=\sum_{j=1}^J w_j\delta_{x_j}$, $w_j\geq0$, $\sum_jw_j=1$, with arbitrary signed $d(x_j)$, the exact endpoints are the minimum and maximum of
$B-\sum_jw_jd(x_j)q_j$ over
\[
 \ell(x_j)\leq q_j\leq u(x_j),\qquad
 |q_j-q_k|\leq L\operatorname{dist}(x_j,x_k)\quad(j,k\leq J).\tag{2}
\]
\end{proposition}
\begin{proof}
Every extension gives a feasible vector. Conversely, the envelope restrictions enforce all Lipschitz inequalities between a grid point and every observed point, and the pairwise constraints enforce those among grid points. Thus the combined function on $S_0\cup\{x_1,\ldots,x_J\}$ is Lipschitz and range bounded. Its lower distance envelope, clipped to $[\epsilon,1-\epsilon]$, extends it to $\mathcal X$. Every feasible vector is consequently attainable. The compact convex feasible set and linear objective prove both endpoint attainment and the interval assertion.
\end{proof}

Choosing a different pointwise envelope according to the sign of $d$ need not yield a Lipschitz function. The pairwise restrictions in (2) are essential; unconstrained pointwise optimization can give unattainable bounds.

\section{A support-gap calculation}
\begin{proposition}[Irreducible width]
Let $\mathcal X=[0,1]$, $S_0=[0,a]$ with $0<a<1$, $p=1/2$ on $S_0$, $d(x)=d_0\geq0$, and $\nu$ uniform on $[0,1]$. If $L(1-a)\leq1/2-\epsilon$, the exact width of (1) is
\[
 d_0 L(1-a)^2.
\]
\end{proposition}
\begin{proof}
On $[0,a]$ the width is zero. On $(a,1]$ the nearest observed point is $a$, and the two envelopes are $1/2\pm L(x-a)$ without clipping under the assumed bound. Integrating $2d_0L(x-a)$ from $a$ to 1 gives the formula.
\end{proof}

\section{An honest region from a finite reference sample}
Continue on $S_0=[0,a]$. Observe $n_0$ independent pairs $(X,A)$ from group 0, with density on this interval at least $f_{\min}>0$. Divide it into $M$ equal bins of width $h=a/M$ and centers $z_j$. Write $N_j$ for bin counts and $\widehat p_j$ for empirical treatment fractions. For $N_j>0$ set
\[
 r_j=\sqrt{\log(2M/\eta)/(2N_j)},\quad
 l_j=\max\{\epsilon,\widehat p_j-r_j-Lh/2\},\quad
 u_j=\min\{1-\epsilon,\widehat p_j+r_j+Lh/2\}.
\]
For an empty bin use $[l_j,u_j]=[\epsilon,1-\epsilon]$. Let $\mathcal P_n$ contain all range-bounded $L$-Lipschitz functions satisfying $l_j\leq p(z_j)\leq u_j$. Project $\mathcal P_n$ through the target functional to obtain $\mathcal C_n$; if it is empty, return the empty set.

\begin{theorem}[Finite-sample coverage and width]
For every admissible reference kernel, $\Prob(D\in\mathcal C_n)\geq1-\eta$. If every bin is nonempty and $r_{\max}=\max_jr_j$, then
\[
 \operatorname{diam}(\mathcal C_n)\leq
 \int |d(x)|\min\{1-2\epsilon,
       2r_{\max}+3Lh+2L\operatorname{dist}(x,S_0)\}\,d\nu(x).\tag{3}
\]
Moreover all $N_j\geq n_0f_{\min}h/2$ except on an event of probability at most $M\exp(-n_0f_{\min}h/8)$. Thus the shrinking part of (3) is $O((\log n_0/n_0)^{1/3})$ for a suitable $h$ of that order, with constants depending on $a,L,f_{\min},\eta$.
\end{theorem}
\begin{proof}
Conditional on the bin counts, each bin's treatment observations are independent Bernoulli draws with mean equal to that bin's conditional mean of $p(X)$. Hoeffding's inequality bounds deviation from this mean by $r_j$ with failure at most $\eta/M$. Lipschitzness bounds the difference between that mean and $p(z_j)$ by $Lh/2$. A union bound therefore places the true kernel in $\mathcal P_n$ with probability at least $1-\eta$, proving coverage.

For any two feasible functions, choose a center whose distance from $x$ is at most $\operatorname{dist}(x,S_0)+h$. Their difference at the center is at most $2r_{\max}+Lh$; adding their two Lipschitz deviations gives the second expression in the minimum of (3). Their common range gives the first. Integrate against $|d|$ to bound the target diameter. Each bin probability is at least $f_{\min}h$. The elementary binomial lower-tail Chernoff bound and a union bound give the count assertion. Substituting this count bound gives $r_{\max}=O(\sqrt{\log n_0/(n_0h)})$; balancing it with $h$ yields the stated shrinking rate. The distance term need not shrink.
\end{proof}

\section{Remaining gap and reference}
The confidence region covers extension uncertainty and reference-sample uncertainty in the declared oracle experiment. Estimating $d,p_1$, allowing higher-order H\"older restrictions and deriving optimal rates for general support geometries require additional work. The residual disparity component is algebraic and is not interpreted as a causal effect of group membership.

Carlos Cinelli, Avi Feller, Guido Imbens, Edward Kennedy, Sara Magliacane and Jose Zubizarreta (2025), \emph{Challenges in Statistics: A Dozen Challenges in Causality and Causal Inference}, arXiv:2508.17099v1, Section 3.4.3. \url{https://arxiv.org/html/2508.17099v1#S3.SS4.SSS3}. This motivates mechanism-based disparity decomposition. The extension and confidence statements above concern the explicitly restricted mathematical experiment.

\end{document}
